% Originally: content/week-11.tex, \section{Gauss' theorem revisited: the case n=1} (Corsin Week 11)


\section{Gauss' theorem revisited: the case \texorpdfstring{$n=1$}{n=1}, and a direct proof
for a square}
\label{sec:gauss_intuition}
% Corsin taught this a week after the divergence theorem, so the source carried a
% \continuedfrom{sec:flux_divergence_theorem} back-pointer. Here the section comes first, as
% motivation, and the pointer would have run forwards; it is gone, and the sentence below
% already \cref's the theorem it previews.
%
% ⚠️ OPEN, raised by the 2026-08-13 audit and deliberately not resolved here. Moving the section
% forward is defensible for the n=1 warm-up, which is motivation and reads perfectly well before
% the statement. It is harder to defend for what follows it: the block below is a genuine
% \begin{proof}[Proof of thm:gauss_divergence for n=2, B=(0,1)^2], and it now sits about ninety
% lines BEFORE the theorem it proves is stated. A reader meets a proof of something they have not
% been told yet. Two repairs are available and each costs something:
%   (a) leave the n=1 warm-up here and move the n=2 proof into sec:flux_divergence_theorem, just
%       under thm:gauss_divergence, where a proof belongs. Splits one source page across two
%       sections, and the "conversely, the FTC can be used to PROVE Gauss" sentence loses the
%       n=1 calculation it answers.
%   (b) relabel the n=2 block as a worked example previewing the theorem rather than proving it,
%       which is closer to what it does in its current position, and keep everything in place.
% (b) is the smaller change and probably the right one. Not taken unilaterally because it alters
% what the block claims to be, and Corsin's own framing is that it proves the theorem.

% Source: Corsin Nick/Class Notes/Week 11.pdf, p. 2
Gauss' theorem (\cref{thm:gauss_divergence}), proved in \cref{sec:flux_divergence_theorem}
below, is, in a precise sense, a higher-dimensional
fundamental theorem of calculus: it trades a derivative integrated over a region for the
original quantity evaluated on the region's boundary. Setting $n=1$ makes the analogy
literal. Let $f \in C^1([a,b])$, viewed as a (one-dimensional) vector field. Since $\divg f =
f'$ in dimension $1$, Gauss' formula becomes
\[\int_a^b \frac{\partial f}{\partial x}\,dx = \int_{[a,b]} \divg f\,d\vol_1
  = \int_{\partial[a,b]} \langle f,\nu\rangle\,d\vol_0
  = f(b) - f(a) \qquad \text{(FTC, almost).}\]
The last step identifies $\int_{\partial[a,b]} \langle f,\nu\rangle\,d\vol_0$ with $f(b)-f(a)$
by treating $\partial[a,b] = \{a,b\}$ as a $0$-dimensional submanifold with outward unit
normals $\nu(a) = -1$, $\nu(b) = +1$, and $\vol_0$ as counting measure.

\begin{ainote}
Corsin flags this himself: \qt{there are several problems with this argument, for example
what a $0$-dim Jordan measure or $0$-dim submanifold is}, and notes that \qt{the definitions
from the lecture can be extended to include this case} without carrying out the extension.
The word \qt{almost} above is his. Not pursued further here; the point of the calculation is
the analogy with the FTC, not a rigorous $n=1$ instance of \cref{thm:gauss_divergence}.
\end{ainote}

% Source: Corsin Nick/Class Notes/Week 11.pdf, p. 3
\begin{proof}[Proof of \cref{thm:gauss_divergence} for $n=2$, $B = (0,1)^2$]
Conversely, the fundamental theorem of calculus can be used to \emph{prove} Gauss' theorem;
we carry this out for the unit square. Let $F : [0,1]^2 \to \mathbb{R}^2$ be a $C^1$ vector
field and $B := (0,1)^2$. The boundary $\partial B$ splits into four edges, with outward unit
normals $-e_1$, $e_1$, $-e_2$, $e_2$ on the left, right, bottom, and top edges respectively.

% Sonnet 5 (Medium) --- FIG-W11-01
\begin{center}
\begin{tikzpicture}[>=Latex, scale=1.6]
  \draw[thick] (0,0) rectangle (1,1);
  \draw[->] (0,0.5) -- (-0.6,0.5) node[left, font=\scriptsize] {$-e_1$};
  \draw[->] (1,0.5) -- (1.6,0.5) node[right, font=\scriptsize] {$e_1$};
  \draw[->] (0.5,0) -- (0.5,-0.6) node[below, font=\scriptsize] {$-e_2$};
  \draw[->] (0.5,1) -- (0.5,1.6) node[above, font=\scriptsize] {$e_2$};
  \node[below left, font=\scriptsize] at (0,0) {$(0,0)$};
  \node[below right, font=\scriptsize] at (1,0) {$(1,0)$};
  \node[above left, font=\scriptsize] at (0,1) {$(0,1)$};
  \node[above right, font=\scriptsize] at (1,1) {$(1,1)$};
\end{tikzpicture}
\end{center}

Parametrizing each edge by $t \mapsto (0,t)$, $(1,t)$, $(t,0)$, $(t,1)$ for $t \in [0,1]$,
\[
\int_{\partial B} \langle F,\nu\rangle\,d\vol_1
  = \int_0^1 \langle F(0,t),-e_1\rangle\,dt + \int_0^1 \langle F(1,t),e_1\rangle\,dt
  + \int_0^1 \langle F(t,0),-e_2\rangle\,dt + \int_0^1 \langle F(t,1),e_2\rangle\,dt.
\]
Grouping the first pair and the second pair,
\[
\int_{\partial B} \langle F,\nu\rangle\,d\vol_1
  = \int_0^1 \big[F_1(1,t) - F_1(0,t)\big]\,dt + \int_0^1 \big[F_2(t,1) - F_2(t,0)\big]\,dt.
\]
By the fundamental theorem of calculus applied to $t \mapsto F_1(s,t)$ along $s$ and to
$s \mapsto F_2(t,s)$ along the second variable,
\[
\int_{\partial B} \langle F,\nu\rangle\,d\vol_1
  = \int_0^1\!\!\left(\int_0^1 \partial_1 F_1(s,t)\,ds\right) dt
  + \int_0^1\!\!\left(\int_0^1 \partial_2 F_2(t,s)\,ds\right) dt,
\]
and by Fubini,
\[
\int_{\partial B} \langle F,\nu\rangle\,d\vol_1
  = \int_B \partial_1 F_1\,d\vol_2 + \int_B \partial_2 F_2\,d\vol_2
  = \int_B \divg F\,d\vol_2. \qedhere
\]
\end{proof}
